% Generated draft: requires mathematical and source review.
\documentclass[11pt]{article}
\usepackage[margin=25mm]{geometry}
\usepackage{fontspec}
\setmainfont{Latin Modern Roman}
\usepackage{amsmath,amssymb,mathtools,unicode-math}
\setmathfont{Latin Modern Math}
\usepackage{enumitem,needspace,xurl,hyperref}
\setlength{\parindent}{0pt}
\setlength{\parskip}{5pt}
\setlength{\emergencystretch}{3em}
\newcommand{\N}{\mathbb{N}}
\newcommand{\Z}{\mathbb{Z}}
\newcommand{\Q}{\mathbb{Q}}
\newcommand{\R}{\mathbb{R}}
\newcommand{\C}{\mathbb{C}}
\newcommand{\F}{\mathbb{F}}
\newcommand{\A}{\mathbb{A}}
\newcommand{\PP}{\mathbb{P}}
\newcommand{\E}{\mathbb{E}}
\newcommand{\eps}{\varepsilon}
\newcommand{\dd}{\,\mathrm{d}}
\begin{document}
\section*{UnsolvedMath ID: 30001222}
\textbf{Source problem code:} catalogue entry
\subsection{Definitions and mathematical statement}
Let \(k\) be a field. All algebras and algebra isomorphisms below are over \(k\), and algebras are understood to be associative with identity. Under the usual presentation convention, a quantum complete intersection is an algebra
\[
X\cong k\langle x_1,\ldots,x_n\rangle\Big/\bigl(x_i^{a_i}\ (1\le i\le n),\ x_ix_j-q_{ij}x_jx_i\ (1\le i<j\le n)\bigr),
\]
where \(n\ge1\), the integers \(a_i\ge2\), and \(q_{ij}\in k^{\times}\). Its dimension is \(\prod_{i=1}^n a_i\). The supplied record does not specify a presentation convention or whether the zero-generator case \(X=k\) is included; the intended convention requires review.

For a \(k\)-algebra \(A\), its center is
\[
Z(A)=\{z\in A:za=az\text{ for every }a\in A\}.
\]
A finite-dimensional algebra \(A\) is local if \(A/J(A)\) is a division algebra, where \(J(A)\) denotes its Jacobson radical. In particular, this definition does not require \(A/J(A)\cong k\). It is symmetric if \(A\cong\operatorname{Hom}_k(A,k)\) as \(A\)-\(A\)-bimodules, with the usual dual bimodule structure.

Here a stable equivalence of Morita type between \(X\) and \(Y\) means that there exist finite-dimensional bimodules \({}_X M_Y\) and \({}_Y N_X\), each projective as a left module and as a right module, such that
\[
M\otimes_Y N\cong X\oplus P,
\qquad
N\otimes_X M\cong Y\oplus Q,
\]
as \(X\)-\(X\)-bimodules and \(Y\)-\(Y\)-bimodules, respectively, where \(P\) is projective over \(X\otimes_kX^{\mathrm{op}}\) and \(Q\) is projective over \(Y\otimes_kY^{\mathrm{op}}\).

\textbf{Question.} For every field \(k\), every quantum complete intersection \(X\) over \(k\), and every finite-dimensional symmetric local \(k\)-algebra \(Y\), suppose that \(X\) and \(Y\) are stably equivalent of Morita type and that
\[
\dim_kY=\dim_kX,
\qquad Z(X)\cong Z(Y)\quad\text{as }k\text{-algebras}.
\]
Must \(X\cong Y\) as \(k\)-algebras?

No restriction on the characteristic or algebraic closure of \(k\), and no additional symmetry hypothesis on \(X\), is stated in the supplied formulation. The supplied dated curation reports a result of Kessar for a nine-dimensional case, but gives insufficient hypotheses to formulate that result precisely here; it is not used to restrict the question.
\subsection{Short English statement}
Must a quantum complete intersection be isomorphic to every finite-dimensional symmetric local algebra of the same dimension and with an isomorphic center to which it is stably equivalent of Morita type? The field is unrestricted in the supplied statement. The reported partial solution and the precise scope of its nine-dimensional case require independent review.
\subsection{Sources}
\begin{itemize}
\item Supplied local UnsolvedMath registry record, permanent ID \(30001222\), source problem code \texttt{catalogue entry}, title: \emph{Rigidity Under Stable Equivalence of Quantum Complete Intersections}. The record belongs to the Oberwolfach Reports Open Problems collection and gives the proposed year as 2009. No individual proposer is identified.
\item Supplied original-report citation: \emph{Representations of Finite Groups} (2009), Oberwolfach Reports, pp.~917--970, DOI \texttt{10.4171/owr/2009/17}. Supplied source URL and literature-source link labelled ``original report'': \url{https://doi.org/10.4171/owr/2009/17}. The supplied background describes Oberwolfach Reports as reports of mathematical workshops and their open-problem discussions.
\item Supplied literature-source link labelled ``Kessar 2012'': \url{https://doi.org/10.1112/jlms/jdr047}. No article title or precise theorem statement is supplied.
\item Literature assessment dated exactly \texttt{2026-08-21}: the supplied status is ``partially solved,'' and the supplied assessment attributes the isomorphism conclusion in a nine-dimensional case to Kessar while describing the general quantum complete intersection question as open. This is supplied dated curation derived from the OWR record, not an independently verified literature assessment. Its time-sensitive status and cited claims require independent verification.
\item The supplied registry labels the statement \texttt{corrected\_verified}; its accompanying verification concerns balanced mathematical delimiters and preservation of hypotheses. That label does not establish mathematical correctness or independent verification of the cited sources.
\end{itemize}
\end{document}
